\section*{Lista de acrónimos}

% \begin{itemize}

% 	\item ESAVI - Evento Supuestamente Atribuible a la Vacunación o Inmunización: término utilizado para describir cualquier evento adverso que se sospecha está relacionado con la administración de una vacuna.

% 	\item VAERS — Vaccine Adverse Event Reporting System: sistema nacional de vigilancia pasiva de eventos adversos en Estados Unidos.
	
% 	\item V-safe — Vaccine Safety Assessment System (CDC): sistema de vigilancia activa mediante encuestas digitales post-vacunación.
	
% 	\item CDC WONDER — Wide-ranging Online Data for Epidemiologic Research: plataforma de acceso público a datos epidemiológicos anonimizados del CDC.
	
% 	\item VSD — Vaccine Safety Datalink: red de vigilancia activa que integra datos de múltiples organizaciones de atención médica para estudios de seguridad vacunal.
	
% 	\item CISA — Clinical Immunization Safety Assessment: programa de evaluación clínica de eventos adversos asociados a la vacunación.
	
% 	\item FDA — Food and Drug Administration: agencia reguladora de medicamentos y alimentos en Estados Unidos.
	
% 	\item CDC — Centers for Disease Control and Prevention: agencia federal responsable del control y prevención de enfermedades.
	
% 	\item AEMPS — Agencia Española de Medicamentos y Productos Sanitarios: organismo regulador de medicamentos en España.
	
% 	\item SEFV-H — Sistema Español de Farmacovigilancia de Medicamentos de Uso Humano: red nacional de notificación de reacciones adversas.
	
% 	\item FEDRA — Farmacovigilancia Española, Datos de Reacciones Adversas: base de datos nacional de eventos adversos.
	
% 	\item OMS — Organización Mundial de la Salud: organismo internacional responsable de la salud pública global.
	
% 	\item VigiBase — base de datos global de farmacovigilancia gestionada por el Centro de Monitoreo de Uppsala en colaboración con la OMS.

% 	\item EudraVigilance — sistema de la Agencia Europea de Medicamentos (EMA) para la gestión y análisis de notificaciones de sospechas de reacciones adversas a medicamentos.
	
% 	\item EMA — European Medicines Agency: agencia reguladora de medicamentos de la Unión Europea.



% 	\item ACID — Atomicity, Consistency, Isolation, Durability: conjunto de propiedades que garantizan la integridad de las transacciones en bases de datos relacionales.
	
% 	\item BASE — Basically Available, Soft state, Eventual consistency: modelo que prioriza la disponibilidad constante del sistema, incluso a costa de la consistencia inmediata, común en bases de datos NoSQL.
	


% \end{itemize}


\newacronym{esavi}{ESAVI}{Evento Supuestamente Atribuible a la Vacunación o Inmunización}
\newacronym{vaers}{VAERS}{Vaccine Adverse Event Reporting System}
\newacronym{vsafe}{v-safe}{Vaccine Safety Assessment System}
\newacronym{cdcwonder}{CDC WONDER}{Wide-ranging Online Data for Epidemiologic Research}
\newacronym{vsd}{VSD}{Vaccine Safety Datalink}
\newacronym{cisa}{CISA}{Clinical Immunization Safety Assessment}
\newacronym{fda}{FDA}{Food and Drug Administration}
\newacronym{cdc}{CDC}{Centers for Disease Control and Prevention}
\newacronym{aemps}{AEMPS}{Agencia Española de Medicamentos y Productos Sanitarios}
\newacronym{sefv}{SEFV-H}{Sistema Español de Farmacovigilancia de Medicamentos de Uso Humano}
\newacronym{fedra}{FEDRA}{Farmacovigilancia Española, Datos de Reacciones Adversas}
\newacronym{oms}{OMS}{Organización Mundial de la Salud}
\newacronym{ema}{EMA}{European Medicines Agency}
\newacronym{acid}{ACID}{Atomicity, Consistency, Isolation, Durability}
\newacronym{base}{BASE}{Basically Available, Soft state, Eventual consistency}


\printglossary[type=\acronymtype, title={Lista de acrónimos}]